\documentclass[conference]{IEEEtran}
\IEEEoverridecommandlockouts
\usepackage{cite}
\usepackage{amsmath,amssymb,amsfonts}
\usepackage{algorithmic}
\usepackage{graphicx}
\usepackage{textcomp}
\usepackage{xcolor}
\usepackage{booktabs}
\usepackage{multirow}

\def\BibTeX{{\rm B\kern-.05em{\sc i\kern-.025em b}\kern-.08em
    T\kern-.1667em\lower.7ex\hbox{E}\kern-.125emX}}

\begin{document}

\title{Evaluating the Robustness of Real-Time Deepfake Detection Under Video-Conferencing Degradation: An Empirical Study}

\author{
\IEEEauthorblockN{1\textsuperscript{st} Ghodake Naman Nitin}
\IEEEauthorblockA{
24BCE5051 \\
SCOPE \\
Vellore Institute of Technology, Chennai \\
naman.ghodake2024@vitstudent.ac.in}
\and
\IEEEauthorblockN{2\textsuperscript{nd} Ekansh Dhanotiya}
\IEEEauthorblockA{
24BCE5127 \\
SCOPE \\
Vellore Institute of Technology, Chennai \\
ekansh.dhanotiya2024@vitstudent.ac.in}
\and
\IEEEauthorblockN{3\textsuperscript{rd} Aditya Gautam}
\IEEEauthorblockA{
24BCE5308 \\
SCOPE \\
Vellore Institute of Technology, Chennai \\
aditya.gautam2024@vitstudent.ac.in}
}

\maketitle

\begin{abstract}
Deepfake detection has been studied extensively for offline, pre-recorded, and public-figure media, yet most published detectors are evaluated under conditions that do not reflect live video-conferencing calls, where compression, bandwidth variance, and other network-induced degradations are the norm rather than the exception. This gap is increasingly consequential, since reported deepfake-enabled video-call fraud has grown sharply over the past year, including cases involving direct financial loss to businesses. In this work we present an empirical robustness study of a fine-tuned Xception-based deepfake detector, evaluated across a systematic matrix of realistic video-call degradations that includes bitrate compression, resolution downscaling, Gaussian blur, and simulated packet loss, applied to a balanced 1,000-video subset of the FaceForensics++ (C23) dataset. We also compare our fine-tuned detector against an off-the-shelf pretrained checkpoint of the same architecture and against accuracy and AUC figures reported for Xception-based detectors in prior published work, in order to contextualize our results. Our detector achieves near-perfect discrimination at baseline (AUC of 0.9997) but degrades unevenly across degradation types. Bitrate compression and simulated packet loss have only a mild effect on detection performance, while resolution downscaling causes a moderate decline (AUC of 0.921 at 240p). The most striking finding is that Gaussian blur is by far the most damaging degradation tested: at a moderate blur level the detector's discriminative power collapses to near-chance performance (AUC of 0.592), even though raw accuracy remains at 0.799 due to a compression of the model's output scores. These results indicate that the fine-grained spatial artifacts on which artifact-based deepfake detectors rely are considerably more vulnerable to loss of spatial detail than to compression or packet loss, a finding with direct implications for how such detectors should be trained and evaluated for real-world, bandwidth- and quality-constrained deployment such as mobile video calls.
\end{abstract}

\begin{IEEEkeywords}
deepfake detection, video conferencing, robustness evaluation, FaceForensics++, Xception, real-time security
\end{IEEEkeywords}

\section{Introduction}

The growing accessibility of deepfake generation tools has turned synthetic media from a research curiosity into an active vector for fraud, impersonation, and disinformation. Early concern in this space centered largely on manipulated offline video: doctored news clips, synthetic celebrity content, and manipulated public-figure footage circulated after the fact. A more operationally dangerous threat has since emerged alongside it, namely real-time deepfake injection into live video calls. Reported cases of deepfake-enabled video-call fraud have risen sharply, growing from a small number of isolated incidents to well over a thousand within a single year in some jurisdictions. High-profile cases have already caused direct financial harm, including at least one instance in which finance staff at a company were defrauded of hundreds of thousands of dollars in a call where every other participant was AI-generated.

Despite this shift in real-world risk, the deepfake detection literature has not kept pace with it. Concrete incidents illustrate the scale of the problem. In February 2024, a finance employee at the engineering firm Arup was reportedly manipulated into transferring approximately 25.5 million U.S. dollars across fifteen separate wire transactions during a video call in which every other participant, including the company's apparent chief financial officer, was an AI-generated deepfake \cite{globalstats2026}. A little over a year later, in March 2025, a finance director at a separate multinational firm in Singapore was reported to have authorized a transfer of approximately 499,000 U.S. dollars after joining what appeared to be a routine video call with senior colleagues, none of whom were real \cite{brightside2025}. These are not isolated anecdotes: one industry threat-intelligence report covering just the first half of 2026 verified 821 distinct deepfake-enabled attacks tied to at least 15,736 individual victims across 3.46 million synthetic files, and a separate analysis of over a billion monitored voice and video interactions reported a 680\% year-over-year increase in deepfake activity \cite{deepstrike2026}. While figures of this kind vary across reporting organizations and should be read as indicative rather than precisely comparable, they consistently point in the same direction, namely that live, interactive impersonation via video and voice has moved from a theoretical risk to an operational one.

The large majority of published detection methods, including the Xception-based baseline established by the widely used FaceForensics++ benchmark \cite{rossler2019faceforensicspp}, are designed and evaluated under offline conditions: high resolution, lightly compressed video, frequently of public figures, with no adversarial network path between the source and the detector. Live video-conferencing calls differ from this setting in several important ways. Video is transmitted through lossy codecs such as H.264, H.265, AV1, and VP9 under variable and often severely constrained bandwidth. Frames may be downscaled to accommodate a poor connection or a low-end device, and packet loss, jitter, and client-side filters such as background blur or beautification further alter the signal before it ever reaches a detector. Each of these factors can strip away exactly the fine-grained pixel-level artifacts that most existing detectors rely on, and the extent to which they do so, and which specific degradation types matter most, has not been systematically characterized in prior work.

This paper addresses that gap through a controlled, empirical robustness study. Rather than proposing a new detection architecture, we take a widely used detector architecture, Xception, fine-tune it on a balanced subset of the FaceForensics++ (C23) dataset, and subject it to a systematic matrix of degradations chosen to approximate real video-call conditions. We evaluate detection performance under each condition using standard metrics (AUC, Equal Error Rate, and threshold accuracy) to identify which degradation types most severely compromise detection and why, and we place our findings in context by comparing against both an untuned baseline checkpoint and figures reported for comparable Xception-based detectors in prior published work.

The contributions of this work are summarized as follows.
\begin{itemize}
\item We construct a reproducible degradation-evaluation pipeline that applies a controlled matrix of bitrate, resolution, blur, and packet-loss-proxy degradations to FaceForensics++ video clips, together with the corresponding fine-tuning, inference, and analysis scripts.
\item We fine-tune an Xception-based detector on clean, undegraded face-cropped frames and evaluate it under each degradation condition at full dataset scale, covering 1,000 videos per condition across four manipulation methods and pristine originals.
\item We identify a previously under-characterized failure mode in which bitrate compression and light simulated packet loss leave detection performance largely intact, while spatial detail loss through resolution downscaling, and especially through Gaussian blur, causes a substantial and in one case near-total collapse in detector discriminability.
\item We show that raw classification accuracy can mask this collapse due to score compression, and argue that AUC and EER based evaluation is necessary to correctly characterize detector robustness under degradation.
\item We compare our fine-tuned detector against an untuned, off-the-shelf pretrained checkpoint of the same architecture, and against figures reported for Xception on FaceForensics++ in prior work, to place our robustness findings in a broader context.
\end{itemize}

\section{Related Work}

\subsection{Offline and Artifact-Based Deepfake Detection}
The majority of deepfake detection research has focused on identifying low-level generation artifacts within individual frames or short clips. The FaceForensics++ benchmark \cite{rossler2019faceforensicspp}, building on the earlier FaceForensics dataset \cite{roessler2018faceforensics}, established a widely used evaluation protocol pairing pristine and manipulated videos across four manipulation methods, namely Deepfakes, Face2Face, FaceSwap, and NeuralTextures, at multiple compression levels. It remains the standard baseline dataset for supervised detector training. CNN-based classifiers, most notably Xception \cite{chollet2017xception} adapted through transfer learning, have served as the reference architecture for this task since the original FaceForensics++ paper. A number of subsequent studies have reported per-method accuracy figures for Xception on FaceForensics++ that we return to in Section~IV as points of comparison; for example, one such study reports binary detection accuracies of 98.98\% on Deepfakes, 97.50\% on Face2Face, 96.48\% on FaceSwap, and 83.00\% on NeuralTextures for an Xception baseline evaluated under a standard, non-degraded protocol \cite{faketransformer2021}. A more recent study reports an in-domain AUC of 0.9805 for Xception trained and evaluated on FaceForensics++ (C23), but only 0.6248 AUC when the same model is evaluated on the unseen FaceShifter manipulation method, illustrating a generalization gap that is a recurring concern in this literature \cite{lightweightfusion2026}. These methods are almost universally evaluated offline, on the dataset's own native compression levels, rather than under the compounding and variable degradations characteristic of live transmission.

\subsection{Real-Time and Video-Call-Specific Detection}
A smaller body of work has begun to specifically target live video-conferencing conditions. Active-probing approaches exploit signals that a live human face reflects correctly but that a synthetic feed cannot replicate under interactive conditions. This includes corneal reflection analysis of on-screen visual challenges \cite{activeprobing2022} and related active-probe techniques \cite{shakingfake2024}, both of which propose challenge-response paradigms as an alternative to passive artifact classification. DF-Captcha \cite{dfcaptcha2022} similarly frames deepfake detection in live calls as an interactive verification problem rather than a purely passive classification task. More recently, industry prototypes such as Fraunhofer SIT's real-time video-conferencing warning system have proposed on-device, privacy-preserving confidence scoring that fuses audio and video signals during a live call \cite{fraunhofer2026}. Taken together, this body of work establishes that real-time detection is an active area of development, but none of it systematically quantifies how a conventional, passively-trained artifact-based detector degrades across the specific technical conditions, namely bitrate, resolution, blur, and packet loss, that a live call actually imposes. That is the gap this paper is intended to fill.

\subsection{Robustness of Detectors Under Image and Video Degradation}
A separate strand of work has studied how deepfake detectors respond to generic image and video degradations, generally without a specific focus on live video-conferencing conditions. A recent large-scale study evaluated five face-swap detectors, including XceptionNet, under JPEG compression, Gaussian noise, motion blur, and reduced-bitrate video encoding on FaceForensics++ and Celeb-DF, and reported that average accuracy fell from 94.7\% on clean data to 67.8\% under distortion, with JPEG compression and motion blur causing the largest drops, of up to 35\%, particularly for lightweight CNN-based detectors such as Xception \cite{robustness2026}. This is closely aligned with our own finding that blur is the most damaging degradation type for an Xception-based detector, and lends independent support to it. A related study evaluating Capsule-Forensics under a range of benign modifications found that Gaussian and mean blur filters reduced accuracy by around 7 to 8 percentage points, a comparatively mild effect relative to what we observe, and attributed the difference in severity across studies partly to the choice of detector architecture and partly to the specific blur kernel or standard deviation used \cite{benignmods2021}. The same study also reported that H.264 encoding at a constant rate factor comparable to our 500 kbps condition reduced accuracy only modestly relative to raw video, consistent with our own observation that bitrate compression is a comparatively mild degradation axis. More broadly, a recent survey-style analysis argued that compression suppresses high-frequency forensic traces while resolution loss distorts the local textures that artifact-based detectors depend on, framing both as manifestations of the same underlying vulnerability to loss of fine spatial detail \cite{loger2026}. Not every reported result points in the same direction, however; at least one study evaluating a diffusion-based forgery detector found that Gaussian blur could, in some cases, mildly improve detection accuracy rather than harm it, suggesting that the direction and magnitude of the blur effect may depend on both the detector architecture and the underlying generation method being detected \cite{diffusionrobust2023}. We revisit these comparisons quantitatively in Section~IV.

\subsection{Research Gap}
Existing real-time-oriented work is largely concerned with proposing new detection paradigms, such as active probing or multi-modal fusion, rather than characterizing the robustness of the conventional artifact-based detectors that remain the most widely deployed and best-understood class of methods. The generic degradation-robustness literature discussed above establishes that compression and blur can harm detector performance in principle, but it is not typically framed around the specific, compounding conditions of a live video call, and it does not isolate resolution downscaling, bitrate compression, blur, and packet loss as separate, directly comparable axes evaluated on the same detector, same dataset, and same evaluation protocol. A systematic, quantitative account of how a single detector's performance decays across these individual, realistic degradation axes, and which axis is most damaging, has not to our knowledge been previously reported at this level of granularity for the video-conferencing setting specifically. This paper aims to provide that account, alongside a grounded comparison against both an untuned baseline and prior literature.

\section{Methodology}

Figure~\ref{fig:architecture} summarizes the overall pipeline used in this study, from dataset construction through fine-tuning, degradation, and evaluation.

\begin{figure}[htbp]
\centering
\includegraphics[width=0.95\linewidth]{architecture_diagram.png}
\caption{Overview of the dataset, fine-tuning, degradation, and evaluation pipeline used in this study.}
\label{fig:architecture}
\end{figure}

\subsection{Dataset}
We use the FaceForensics++ dataset at C23 (lightly compressed) quality \cite{rossler2019faceforensicspp}, which provides pristine (original) video clips alongside four manipulation methods: Deepfakes, Face2Face, FaceSwap, and NeuralTextures. From this dataset we sampled a balanced subset of 200 clips per class (original and each of the four manipulation methods), giving 1,000 source videos in total. Because the real class is drawn from a single source category while the fake class spans four manipulation methods, the resulting subset has a real-to-fake ratio of 1:4, that is, 200 real videos against 800 fake videos. This class imbalance is accounted for in our evaluation protocol, described in Section~III-F.

\subsection{Baseline Detector Architecture}
We use an Xception-based binary classifier, following the transfer-learning protocol established as the FaceForensics++ baseline: an ImageNet-pretrained Xception backbone with its final fully-connected layer replaced by a dropout layer ($p=0.5$) followed by a two-class linear classification head. Faces are localized per frame using a HOG-based frontal face detector, with the detected bounding box enlarged by a fixed scale factor of 1.3 before cropping, so that some facial context beyond the tightest bounding box is preserved. Cropped faces are resized to $299\times299$ and normalized to $[-1, 1]$ before being passed to the network.

\subsection{Fine-Tuning Protocol}
An initial evaluation using a publicly available FaceForensics++-pretrained Xception checkpoint revealed unreliable, high-variance predictions on our sampled clean videos, including frequent misclassification of pristine content as manipulated even in the absence of any applied degradation. To obtain a reliable baseline detector for the robustness study, we fine-tuned the same Xception architecture on our own data rather than relying on this checkpoint directly. Ten evenly spaced, face-cropped frames were extracted from each of the 1,000 clean, undegraded source videos. The resulting frame set was split at the video level, not the frame level, into 85\% training and 15\% validation partitions, so that frames from the same source video never appear in both splits. The network was fine-tuned from the pretrained checkpoint for 5 epochs using the Adam optimizer with a learning rate of $1\times10^{-3}$, $\beta_1=0.9$, $\beta_2=0.999$, a batch size of 8, and cross-entropy loss. The checkpoint corresponding to the epoch with the highest validation accuracy was retained as the final baseline detector, reaching 91.46\% validation accuracy at the first epoch, after which validation accuracy declined slightly across subsequent epochs while training accuracy continued to rise, a pattern consistent with mild overfitting on a comparatively small fine-tuning set.

\subsection{Degradation Matrix}
To approximate the conditions encountered in live video-conferencing calls, we applied a systematic matrix of degradations to every clean source video using FFmpeg, in addition to an unmodified baseline condition. Table~\ref{tab:degradation_matrix} summarizes the conditions evaluated.

\begin{table}[htbp]
\caption{Degradation matrix applied to each source video}
\label{tab:degradation_matrix}
\centering
\begin{tabular}{@{}ll@{}}
\toprule
\textbf{Degradation type} & \textbf{Conditions} \\
\midrule
Baseline & Unmodified (container copy) \\
Bitrate (H.264) & 500, 1000, 2000, 5000 kbps \\
Resolution & 240p, 360p, 480p, 720p \\
Gaussian blur & $\sigma = 1.5$, $\sigma = 3.0$ \\
Simulated packet loss & 1\% (noise-filter proxy) \\
\bottomrule
\end{tabular}
\end{table}

Simulated packet loss was approximated using a per-frame noise-injection filter, since genuine packet-loss simulation requires network-level emulation that falls outside the scope of an offline, file-based degradation pipeline. During pipeline development we found that this noise-based proxy at higher intensities, specifically 5\% and 10\%, produced disproportionate increases in encoded file size under constant-quality encoding, which indicated that it was introducing artifacts inconsistent with how real packet loss actually manifests visually. We therefore restricted the packet-loss condition to a single, mild 1\% setting and treat it as a coarse reference point rather than a faithful reconstruction of network loss.

\subsection{Implementation and Reproducibility}
All experiments were implemented in Python using PyTorch for model training and inference, dlib for face detection, and FFmpeg for the degradation pipeline. Fine-tuning and inference were both carried out on a single consumer-grade GPU (an NVIDIA RTX 3050 with 6 GB of video memory), with the degradation matrix itself generated on CPU, since FFmpeg's encoding filters do not require GPU acceleration. This choice was deliberate: since one motivation for this study is understanding detector behavior under realistic, resource-constrained conditions, we considered it appropriate to also constrain our own experimental hardware to a widely available consumer configuration rather than a data-center-grade setup. Informal timing measurements taken during pipeline development indicated that the dominant cost per video was the per-frame, CPU-bound HOG face detector rather than the Xception network itself, meaning that inference throughput in this pipeline is more a function of face-localization speed than of the classifier's own forward-pass cost. This observation is relevant to real-time deployment discussed in Section~IV-J, since it suggests that optimizing or replacing the face-localization stage, rather than the classifier, may be the more direct route to reducing end-to-end latency in a live system. All dataset sampling, degradation, fine-tuning, and evaluation scripts used in this study operate on the publicly available FaceForensics++ dataset and open-source tooling only, and the balanced 1,000-video subset, degradation matrix, and evaluation protocol described above are fully specified so that the pipeline can be reproduced or extended by other researchers.

\subsection{Evaluation Protocol}
For each of the twelve conditions, that is the baseline plus eleven degraded variants, all 1,000 source videos were processed and evaluated. For each video, up to ten evenly spaced frames were sampled. For each sampled frame, the largest detected face was cropped and classified, giving a per-frame fake-probability score from the network's softmax output. A single video-level fake-probability score was computed as the mean of its per-frame scores. Using these video-level scores together with the known ground-truth label for each video, we computed, per condition: Area Under the ROC Curve (AUC), Equal Error Rate (EER, the point at which false-positive and false-negative rates are equal), and classification accuracy at a fixed decision threshold of 0.5. AUC and EER are threshold-independent and rank-based, and are therefore robust to the class imbalance described in Section~III-A. Raw accuracy is reported alongside them but interpreted with this imbalance in mind.

\section{Results and Discussion}

\subsection{Overall Performance Across Degradation Types}
Table~\ref{tab:results} reports AUC, EER, and accuracy for all twelve evaluated conditions, each computed over the full 1,000-video subset.

\begin{table}[htbp]
\caption{Detector performance across degradation conditions (1,000 videos per condition)}
\label{tab:results}
\centering
\begin{tabular}{@{}llccc@{}}
\toprule
\textbf{Condition} & \textbf{Value} & \textbf{AUC} & \textbf{EER} & \textbf{Acc.} \\
\midrule
Baseline & original & 0.9997 & 0.0075 & 0.995 \\
\midrule
\multirow{4}{*}{Bitrate} & 500 kbps & 0.9743 & 0.0650 & 0.945 \\
 & 1000 kbps & 0.9963 & 0.0150 & 0.979 \\
 & 2000 kbps & 0.9991 & 0.0100 & 0.992 \\
 & 5000 kbps & 0.9991 & 0.0150 & 0.996 \\
\midrule
\multirow{4}{*}{Resolution} & 240p & 0.9207 & 0.1600 & 0.875 \\
 & 360p & 0.9829 & 0.0650 & 0.946 \\
 & 480p & 0.9933 & 0.0500 & 0.966 \\
 & 720p & 0.9960 & 0.0300 & 0.981 \\
\midrule
\multirow{2}{*}{Blur} & $\sigma=1.5$ & 0.9010 & 0.1975 & 0.821 \\
 & $\sigma=3.0$ & \textbf{0.5917} & \textbf{0.4525} & 0.799 \\
\midrule
Packet loss & 1\% & 0.9857 & 0.0550 & 0.890 \\
\bottomrule
\end{tabular}
\end{table}

At baseline, the fine-tuned detector performs close to perfectly, with an AUC of 0.9997 and an EER of 0.0075, confirming that the fine-tuning procedure described in Section~III-C produced a well-calibrated detector on in-distribution, undegraded data. Performance under degradation, however, varies considerably by degradation type rather than declining uniformly with a single notion of signal quality, a pattern that becomes visually clear once the twelve conditions are placed side by side in Figures~\ref{fig:auc_all} and~\ref{fig:eer_all} below.

Figure~\ref{fig:auc_all} presents AUC across all twelve conditions, with bars colored by degradation family. Read left to right, the chart makes the central pattern of this study immediately visible. The blue bitrate bars and the green packet-loss bar remain close to the gray baseline bar across their full range, the orange resolution bars decline gradually and moderately as resolution drops, and the two red blur bars behave differently from every other family, with the second of the two falling sharply below every other condition tested, including the most severe bitrate and resolution conditions. The dotted horizontal line marks the AUC value expected of random guessing, and only the strong blur condition comes close to it, which is itself a striking result given that the other eleven conditions all remain comfortably in the upper half of the plotted range.

\begin{figure}[htbp]
\centering
\includegraphics[width=0.95\linewidth]{auc_all_conditions.png}
\caption{AUC across all twelve evaluated conditions, grouped by degradation family. The dotted line marks the AUC expected of random guessing.}
\label{fig:auc_all}
\end{figure}

Table~\ref{tab:results} and Figure~\ref{fig:auc_all} together establish the headline pattern of this study, but AUC alone does not fully convey how a detector actually behaves at a fixed decision boundary in deployment, which is why we also report Equal Error Rate for every condition.

Figure~\ref{fig:eer_all} presents the same twelve conditions in terms of Equal Error Rate, which in a loose sense mirrors the AUC view from the opposite direction: whereas the AUC chart shows discriminative power falling, the EER chart shows the rate of confusion between the two classes rising by a corresponding amount. The strong blur condition again stands apart from the rest of the matrix, with an EER close to 0.45, meaning that at the operating point where false-positive and false-negative rates are balanced, the detector is wrong on nearly half of all decisions, a result that is difficult to distinguish in practice from not running a detector at all. The mild blur condition and the most severe resolution condition form a second, considerably less severe tier, with EER values in the range of 0.16 to 0.20, while every bitrate condition and the packet-loss condition remain below an EER of 0.07, reinforcing the same ordering of degradation severity already visible in the AUC chart.

\begin{figure}[htbp]
\centering
\includegraphics[width=0.95\linewidth]{eer_all_conditions.png}
\caption{Equal Error Rate across all twelve evaluated conditions, grouped by degradation family.}
\label{fig:eer_all}
\end{figure}

\subsection{Bitrate Compression and Packet Loss Are Comparatively Mild}
Even at the most aggressive bitrate setting tested, 500 kbps, AUC remains high at 0.9743, with a modest EER of 0.065. Performance is effectively back at baseline by 2000 kbps and above, shown in Figure~\ref{fig:bitrate}. The simulated 1\% packet-loss condition likewise has only a mild effect on AUC, at 0.9857, though its accuracy of 0.890 drops more noticeably, a discrepancy we return to in Section~IV-D. This pattern suggests that the discriminative artifacts the detector relies on are, to a meaningful extent, robust to the information loss introduced by standard compression, a plausible outcome given that H.264 compression at these bitrates, while visually lossy, largely preserves macro-level facial structure at the original spatial resolution.

\begin{figure}[htbp]
\centering
\includegraphics[width=0.95\linewidth]{accuracy_vs_bitrate.png}
\caption{Detector AUC and accuracy across the bitrate degradation sweep.}
\label{fig:bitrate}
\end{figure}

\subsection{Resolution Downscaling Causes Moderate but Clear Degradation}
Resolution reduction shows a clearer, monotonic decline in performance as resolution decreases (Figure~\ref{fig:resolution}). AUC falls from 0.9960 at 720p to 0.9207 at 240p, with EER rising from 0.030 to 0.160 over the same range. Unlike bitrate compression, which preserves spatial resolution while discarding high-frequency detail through quantization, downscaling directly destroys the fine-grained spatial detail in which subtle manipulation artifacts, such as blending boundaries and texture inconsistencies, tend to reside. This is an early indication that loss of spatial detail specifically, rather than loss of signal fidelity in a general sense, is the more consequential axis of degradation for this class of detector.

\begin{figure}[htbp]
\centering
\includegraphics[width=0.95\linewidth]{accuracy_vs_resolution.png}
\caption{Detector AUC and accuracy across the resolution degradation sweep.}
\label{fig:resolution}
\end{figure}

\subsection{Gaussian Blur Is the Dominant Failure Mode}
The clearest and most severe result in this study concerns Gaussian blur, shown in Figure~\ref{fig:blur}. At a moderate blur level, $\sigma=1.5$, AUC already falls to 0.9010, comparable to or worse than the most severe bitrate and resolution conditions tested. At a stronger blur level, $\sigma=3.0$, AUC collapses to 0.5917, barely above the 0.5 value expected of random guessing, meaning the detector has lost almost all ability to rank real videos above fake ones under this condition.

Notably, accuracy at $\sigma=3.0$ remains at 0.799, a figure that in isolation might appear only moderately degraded. This divergence between a near-random AUC and a superficially reasonable accuracy is a signature of score compression: under severe blur, the network's output probabilities cluster tightly around a narrow range, such that a fixed 0.5 decision threshold happens to produce a plausible-looking accuracy figure even though the underlying score distribution carries almost no discriminative information. This finding underscores a broader methodological point, namely that accuracy alone, evaluated at a single fixed threshold, can substantially understate the true collapse of a detector's discriminative capability under degradation. Threshold-independent metrics such as AUC and EER are necessary to correctly characterize robustness in this setting.

\begin{figure}[htbp]
\centering
\includegraphics[width=0.95\linewidth]{accuracy_vs_blur_packetloss.png}
\caption{Detector AUC and accuracy under Gaussian blur and simulated packet loss. Note the divergence between AUC and accuracy at $\sigma=3.0$.}
\label{fig:blur}
\end{figure}

The severity of the blur result relative to resolution downscaling is notable, given that both operations reduce spatial detail. We attribute the larger effect of blur to the fact that Gaussian blur removes high-frequency spatial information at the original resolution, directly attacking the fine-grained pixel-level inconsistencies, such as blending seams and texture discontinuities, that artifact-based detectors are trained to key on, without the partially compensating effect of a smaller but still-sharp image that resolution downscaling provides.

\subsection{Relation to the Broader Degradation-Robustness Literature}
Our finding that blur is the dominant failure mode is consistent with, and to our knowledge is one of the more severe instances of, a pattern that has begun to appear elsewhere in the literature. A recent evaluation of five face-swap detectors, including XceptionNet, under JPEG compression, Gaussian noise, motion blur, and reduced-bitrate encoding reported that average accuracy across all five models and all four distortion types fell from 94.7\% on clean data to 67.8\% on distorted data, with compression and motion blur singled out as causing the largest individual drops, of up to 35\%, particularly for lightweight CNN-based detectors such as Xception \cite{robustness2026}. Our own accuracy figures move in the same direction: relative to our 99.5\% baseline accuracy, accuracy falls to 82.1\% under mild blur and 79.9\% under strong blur, drops of roughly 17 and 20 percentage points respectively, while our AUC collapses far more severely, from 0.9997 to 0.5917, a gap between the accuracy-based and AUC-based view of the same condition that is itself one of this paper's central observations, discussed above in Section~IV-D.

By contrast, a study of Capsule-Forensics under similar benign modifications reported a considerably smaller blur effect, with Gaussian and mean blur filters reducing accuracy by only around 7 to 8 percentage points, and attributed cross-study differences in severity partly to detector architecture and partly to the specific blur kernel size or standard deviation applied \cite{benignmods2021}. Our own use of a comparatively strong Gaussian blur, $\sigma=3.0$, at the more severe end of what is typically tested, likely accounts for at least part of the difference in magnitude we observe relative to that study. The same work reported that H.264 compression at a quality setting broadly comparable to our most aggressive 500 kbps condition produced only a modest reduction in accuracy relative to raw video, which agrees with our own observation that bitrate compression is a comparatively mild degradation axis for this class of detector. A broader discussion of degradation robustness in recent detector design work frames both compression and resolution loss as manifestations of the same underlying vulnerability, namely that artifact-based detectors depend on fine spatial and high-frequency detail that either operation can remove \cite{loger2026}, which is consistent with the mechanism we propose in Section~IV-D for why blur is more damaging than resolution downscaling at comparable severity.

Not every result in this space points the same way. One study of a diffusion-based forgery detector found that Gaussian blur could, in some cases, mildly improve rather than harm detection accuracy, which suggests that the direction and severity of the blur effect may be specific to both the detector architecture and the family of generation method being detected, rather than a universal property of blur itself \cite{diffusionrobust2023}. Our results should therefore be read as characterizing this particular Xception-based detector trained on FaceForensics++-style manipulations, rather than as a claim that blur is universally the worst-case degradation for every deepfake detector or generation method.

\subsection{Comparison with an Off-the-Shelf Pretrained Baseline}
As noted in Section~III-C, our decision to fine-tune the Xception architecture on our own data was motivated by unreliable behavior observed when using a publicly available, off-the-shelf FaceForensics++-pretrained checkpoint of the same architecture. To make this comparison concrete, Figure~\ref{fig:pretrained_finetuned} shows the predicted fake-probability of both the pretrained checkpoint and our fine-tuned detector on the same five held-out, undegraded real videos, evaluated under identical baseline conditions.

\begin{figure}[htbp]
\centering
\includegraphics[width=0.95\linewidth]{pretrained_vs_finetuned.png}
\caption{Predicted fake probability on five held-out real videos, comparing the off-the-shelf pretrained checkpoint against our fine-tuned detector. Lower is correct, since these are all pristine videos.}
\label{fig:pretrained_finetuned}
\end{figure}

The off-the-shelf checkpoint assigns a fake probability above 0.94 to three of the five videos and above 0.99 to two of those three, despite all five being genuine, unmanipulated footage. Our fine-tuned detector assigns fake probabilities between 0.17 and 0.55 to the same five videos, a substantial and consistent improvement. This comparison is based on a small sample of five videos and should be read as an illustrative, motivating observation rather than a statistically powered result, but it is consistent with, and helps explain, the large-scale baseline accuracy of 99.5\% we report for the fine-tuned detector in Table~\ref{tab:results}.

\subsection{Comparison with Literature-Reported Baselines}
Table~\ref{tab:literature} places our baseline result in the context of Xception-based accuracy and AUC figures reported for FaceForensics++ in prior published work. These figures are not directly comparable to ours in a strict sense, since they were obtained under different evaluation protocols, different data subsets, and in most cases frame-level rather than video-level aggregation, and none of them involve the degradation matrix used in this study. They are included only as approximate reference points for how our baseline detector's performance sits relative to previously reported results for the same underlying architecture.

\begin{table}[htbp]
\caption{Baseline (undegraded) performance in context of prior reported work}
\label{tab:literature}
\centering
\begin{tabular}{@{}lcc@{}}
\toprule
\textbf{Source} & \textbf{Metric} & \textbf{Value} \\
\midrule
Xception, per-method accuracy \cite{faketransformer2021} & Accuracy & 83.0--99.0\% \\
Xception, in-domain (C23) \cite{lightweightfusion2026} & AUC & 0.9805 \\
Xception, cross-manipulation \cite{lightweightfusion2026} & AUC & 0.6248 \\
Public checkpoint, in-domain \cite{mumairab2026} & AUC & 0.9944 \\
Public checkpoint, cross-manipulation \cite{mumairab2026} & AUC & 0.70 \\
\textbf{Ours, baseline (undegraded)} & \textbf{AUC} & \textbf{0.9997} \\
\bottomrule
\end{tabular}
\end{table}

Our baseline AUC is at or above the range of these previously reported figures, which is consistent with our fine-tuning having produced a competently calibrated detector, though we emphasize again that the comparison is only approximate given the protocol differences involved. What is more striking is the recurring cross-manipulation and cross-condition generalization gap visible across every row of Table~\ref{tab:literature}: even detectors that perform well in-domain lose a substantial fraction of their AUC when the evaluation distribution shifts, whether that shift is a change in manipulation method, as in \cite{lightweightfusion2026,mumairab2026}, or, as we show in Sections~IV-B through IV-D, a change in video-call degradation conditions. This reinforces the view that in-domain accuracy figures alone are an incomplete measure of a detector's practical reliability.

\subsection{Qualitative Examples}
To illustrate the detector's behavior directly, rather than only through aggregate metrics, this section presents representative frames processed through the detector's own inference script, with the predicted class probabilities overlaid directly on the frame.

Figure~\ref{fig:screenshot_real1} shows a frame from the first of two genuine, unmanipulated sample videos. The green bounding box marks the face region located by the dlib detector before cropping, and the overlaid text reports the model's raw two-class softmax output, ordered as real followed by fake, alongside the resulting label. In this frame the detector assigns essentially all of its predicted probability mass to the real class, correctly identifying the footage as genuine with high confidence.

\begin{figure}[htbp]
\centering
\includegraphics[width=0.85\linewidth]{screenshot_real_1.png}
\caption{Detector output on a genuine, unmanipulated sample video, correctly predicting real with a class probability of 1.00 real / 0.00 fake.}
\label{fig:screenshot_real1}
\end{figure}

Figure~\ref{fig:screenshot_real2} shows the same evaluation applied to a second genuine video, drawn from a different source and studio setting than the first, with a different subject, framing, and background. The detector again correctly predicts real with a class probability of 1.00 real and 0.00 fake. While this is only a qualitative pair of examples rather than a systematic subgroup analysis, the fact that both are classified correctly and with similarly high confidence, despite differing in subject and visual setting, is at least mildly reassuring evidence against the model having overfit to some narrow, superficial visual cue specific to a single video or presenter.

\begin{figure}[htbp]
\centering
\includegraphics[width=0.85\linewidth]{screenshot_real_2.png}
\caption{Detector output on a second genuine, unmanipulated video, again correctly predicting real with a class probability of 1.00 real / 0.00 fake.}
\label{fig:screenshot_real2}
\end{figure}

Figure~\ref{fig:screenshot_fake} is included as a placeholder for the corresponding example on a manipulated video. We did not capture a screenshot of the detector's output on a fake video during development in a form suitable for direct inclusion here, and in keeping with our commitment to reporting only verified figures and genuine artifacts throughout this paper, we have chosen to mark this explicitly as a placeholder rather than substitute a fabricated or reconstructed image. Before final submission, this figure should be replaced with an equivalent screenshot captured from the same inference script run on one of the Deepfakes, Face2Face, FaceSwap, or NeuralTextures videos in our subset, together with its exact overlaid probability values reported in the caption, following the same convention used for Figures~\ref{fig:screenshot_real1} and~\ref{fig:screenshot_real2}.

\begin{figure}[htbp]
\centering
\includegraphics[width=0.85\linewidth]{screenshot_fake_PLACEHOLDER.png}
\caption{Placeholder for a detector output on a manipulated (fake) video, to be replaced with a screenshot from the authors' own inference run before submission.}
\label{fig:screenshot_fake}
\end{figure}

\subsection{Implications for Real-World Deployment}
These results carry direct implications for deploying artifact-based deepfake detectors against live video-call fraud. The video-call scenarios in which fraud is most commonly reported, namely mobile connections, low-bandwidth environments, and calls using background-blur or beautification filters, are precisely the scenarios in which spatial detail loss, whether through native low resolution or client-side blur filters, is most likely, rather than scenarios dominated by bitrate starvation alone. A detector validated only against compression-based degradation may therefore substantially overstate its real-world reliability. This suggests that resolution- and blur-aware training augmentation, rather than compression-aware augmentation alone, should be prioritized when developing detectors intended for live deployment, and that evaluation protocols for such detectors should explicitly separate spatial-detail-loss conditions from bitrate-driven conditions rather than treating degradation as a single axis. This recommendation is also consistent with the broader degradation-robustness literature discussed in Section~IV-E, which has begun to frame compression and resolution loss as two expressions of a single underlying vulnerability to loss of fine spatial detail rather than as unrelated failure modes \cite{loger2026}; our results add to that picture a direct, controlled comparison showing that among the specific degradations we tested, blur is markedly more damaging than either bitrate compression or the resolution downscaling with which it is sometimes grouped.

\subsection{Computational Considerations for Real-Time Deployment}
Robustness to degradation is only one requirement for a detector intended to run against a live video call; the other is that it must run fast enough to be useful before the call has ended. As noted in Section~III-E, informal timing observations from our own pipeline indicated that face localization, rather than the Xception classifier itself, was the dominant cost per frame when using a CPU-bound HOG-based detector. This observation is consistent with the broader design direction taken by industry prototypes in this space; the Fraunhofer SIT system discussed in Section~II-B, for instance, is explicitly built around on-device processing of both audio and video signals for privacy and latency reasons rather than routing a call through a remote server \cite{fraunhofer2026}. Our results suggest that a similar principle applies at a finer grain within the detection pipeline itself: if a lighter-weight or GPU-accelerated face localization stage were substituted for the HOG-based detector used here, the effective per-frame cost of running the classifier would likely be dominated by the Xception forward pass rather than by face detection, which is a comparatively well-understood and already heavily optimized workload on modern consumer and mobile GPUs. We did not conduct a formal, controlled latency benchmark as part of this study, and we present this only as a qualitative observation from development rather than as a quantitative result, but we believe it is a reasonable and practically useful direction to flag for future work aimed specifically at real-time deployment, alongside the robustness-focused recommendations above.

\section{Limitations}
Several limitations should be kept in mind when interpreting these results. First, an initial off-the-shelf pretrained checkpoint was found to be unreliable on our data, which is why we fine-tuned our own baseline detector as described in Section~III-C. This means our results characterize the robustness of a detector fine-tuned on a relatively modest, single-dataset sample, rather than a large-scale, extensively validated production model, and the pretrained-versus-fine-tuned comparison in Section~IV-F is based on only five videos and should be read as illustrative. Second, our simulated packet-loss condition is a noise-based proxy rather than a true network-level emulation of packet loss, and is limited to a single mild setting for the reasons discussed in Section~III-D; the effect of more severe packet loss on detection performance remains an open question. Third, this study evaluates a single detector architecture, Xception, and we do not establish whether the blur-dominant degradation pattern observed here generalizes to other detector families, such as those based on frequency-domain analysis or vision transformers. Fourth, the literature comparison in Section~IV-G should be read as an approximate, contextual reference rather than a controlled head-to-head comparison, since the cited figures were obtained under different protocols and datasets. Finally, evaluation was conducted on a 1,000-video subset of FaceForensics++ rather than the full dataset and on a single source dataset, so cross-dataset generalization was not assessed here.

\section{Conclusion and Future Work}
This paper presented a systematic empirical study of how a fine-tuned Xception-based deepfake detector's performance degrades under realistic video-conferencing conditions. Contrary to the common framing of compression robustness as the primary challenge for real-time deepfake detection, our results indicate that spatial detail loss, and Gaussian blur in particular, is a substantially more severe failure mode than bitrate compression or simulated packet loss, in one condition reducing detector discriminability to near-chance levels while leaving raw accuracy deceptively stable. We further showed that fine-tuning on in-distribution data produced a considerably more reliable baseline than an off-the-shelf pretrained checkpoint of the same architecture, and that our resulting baseline performance is consistent with, and in some cases favorable relative to, figures reported for comparable detectors in prior published work. Our central blur finding is independently corroborated in direction, though not in exact magnitude, by at least one other recent robustness study of Xception-based face-swap detectors, which strengthens our confidence that this is a genuine property of the architecture rather than an artifact of our particular pipeline, while a contrasting result reported for a diffusion-based detector suggests this vulnerability may not generalize to every detector and generation method equally. These findings motivate several directions for future work: training-time augmentation with blur and resolution degradation to improve robustness at the source; evaluation of whether the blur-dominant vulnerability generalizes across other detector architectures, including frequency-domain and multi-modal audio-video fusion approaches; and replication at larger scale and against true network-level packet-loss emulation to validate and extend the degradation matrix used here.

\begin{thebibliography}{00}

\bibitem{roessler2018faceforensics}
A. R\"ossler, D. Cozzolino, L. Verdoliva, C. Riess, J. Thies, and M. Nie{\ss}ner, ``FaceForensics: A Large-Scale Video Dataset for Forgery Detection in Human Faces,'' \emph{arXiv preprint arXiv:1803.09179}, 2018.

\bibitem{rossler2019faceforensicspp}
A. R\"ossler, D. Cozzolino, L. Verdoliva, C. Riess, J. Thies, and M. Nie{\ss}ner, ``FaceForensics++: Learning to Detect Manipulated Facial Images,'' in \emph{Proc. IEEE/CVF Int. Conf. Computer Vision (ICCV)}, 2019.

\bibitem{chollet2017xception}
F. Chollet, ``Xception: Deep Learning with Depthwise Separable Convolutions,'' in \emph{Proc. IEEE Conf. Computer Vision and Pattern Recognition (CVPR)}, 2017.

\bibitem{activeprobing2022}
``Detection of Real-time DeepFakes in Video Conferencing with Active Probing and Corneal Reflection,'' \emph{arXiv preprint arXiv:2210.14153}, 2022.

\bibitem{shakingfake2024}
``Shaking the Fake: Detecting Deepfake Videos in Real Time via Active Probes,'' \emph{arXiv preprint arXiv:2409.10889}, 2024.

\bibitem{dfcaptcha2022}
``DF-Captcha: A Deepfake Captcha for Preventing Fake Calls,'' \emph{arXiv preprint arXiv:2208.08524}, 2022.

\bibitem{fraunhofer2026}
Fraunhofer Institute for Secure Information Technology (SIT), ``A Real-Time Warning System for Fake Videoconferences,'' Fraunhofer SIT Press Release, Aug. 2026.

\bibitem{faketransformer2021}
``FakeTransformer: Exposing Face Forgery From Spatial-Temporal Representation Modeled By Facial Pixel Variations,'' \emph{arXiv preprint arXiv:2111.07601}, 2021.

\bibitem{lightweightfusion2026}
``Lightweight Complementary-Cue Fusion for Robust Video Face Forgery Detection,'' \emph{arXiv preprint arXiv:2605.29092}, 2026.

\bibitem{mumairab2026}
MUmairAB, ``deepfake-detection-ff-cn-transformer,'' Hugging Face model repository. [Online]. Available: https://huggingface.co/MUmairAB/deepfake-detection-ff-cn-transformer

\bibitem{robustness2026}
``Evaluating the real-world robustness of face-swap detection models under compression and noise,'' \emph{Frontiers in Artificial Intelligence}, 2026.

\bibitem{benignmods2021}
``Impact of Benign Modifications on Discriminative Performance of Deepfake Detectors,'' \emph{arXiv preprint arXiv:2111.07468}, 2021.

\bibitem{loger2026}
``LOGER: Local-Global Ensemble for Robust Deepfake Detection in the Wild,'' \emph{arXiv preprint arXiv:2604.03558}, 2026.

\bibitem{diffusionrobust2023}
``Robustness and Generalizability of Deepfake Detection: A Study with Diffusion Models,'' \emph{arXiv preprint arXiv:2309.02218}, 2023.

\bibitem{globalstats2026}
``Deepfake Statistics 2026: Cases, Victims and Key Facts,'' The Global Statistics, 2026. [Online]. Available: https://www.theglobalstatistics.com/deepfake-statistics/

\bibitem{brightside2025}
``Deepfake CEO Fraud: The Voice Cloning Threat Facing CFOs,'' Brightside AI, 2025. [Online]. Available: https://www.brside.com/blog/deepfake-ceo-fraud-50m-voice-cloning-threat-cfos

\bibitem{deepstrike2026}
``Deepfake Statistics 2026: Fraud, Identity and Detection,'' DeepStrike, 2026. [Online]. Available: https://deepstrike.io/blog/deepfake-statistics-2025

\end{thebibliography}

\end{document}
